% Appendix: judge prompts. Generated by code/llm_judge/make_prompt_appendix.py

\subsection*{Behaviour-judging prompts}

Counts of backtracking and verification come from an LLM judge
(\texttt{gemini-2.5-flash}, temperature 0) given one whole reasoning trace and one
prompt per behaviour. The prompts are reproduced verbatim from
\citet{gandhi2025cognitive}'s \texttt{pretraining\_analysis/prompts/}, with a
single edit: their opening line, ``You will be provided with text from the
internet.'', is replaced by the line marked \textbf{[edited]} below, because we
apply them to reasoning traces rather than web text. Nothing else is changed, and
there is no system message --- their pipeline sends the formatted template as the
only user message. The trace is substituted for \texttt{\{response\}}.

\paragraph{Backtracking.}
\begin{quote}\small\itshape
\# Task Description\\
\textbf{[edited]} You will be provided with the reasoning trace of a language model solving a competition mathematics problem.\\
Evaluate whether the text contains any backtracking behavior, where the writer realizes a path won't work and explicitly goes back to try a different approach. An example of backtracking is: "Let me try again", "Wait", "I made a mistake", or "we need to try a different sequence of operations". We want to mark instances where the writer abandons a thought and backtracks to a previous computation.\\
\medskip
Backtracking in mathematics might look like:\\
- "I started with the wrong formula. Let's use integration by parts instead."\\
- "This approach leads to a contradiction. Going back to the original equation..."\\
- "I see the error in my calculation. Let's recalculate using..."\\
- "This algebraic manipulation isn't simplifying as expected. Let's try factoring differently."\\
\medskip
Count the number of distinct backtracking instances and provide the count between the tags <count> </count>. If the writer does not backtrack, please provide a count of 0 as <count>0</count>.\\
\medskip
\# Task Format\\
Format your response in markdown as follows:\\
\medskip
\#\# Thoughts\\
[Brief description describing what behavior was noticed and where backtracking occurred]\\
\medskip
\#\# Does backtrack?\\
[yes/no]\\
\medskip
\#\# Number of backtrack steps\\
<count> [1/2/...] </count>\\
\medskip
\# Task to evaluate for backtracking\\
\{response\}\\
\medskip
\# Response\\
\end{quote}

\paragraph{Verification.}
\begin{quote}\small\itshape
\# Task Description\\
\textbf{[edited]} You will be provided with the reasoning trace of a language model solving a competition mathematics problem.\\
Evaluate whether the text contains any verification steps. We want to mark instances where the writer explicitly checks their own work, such as by comparing the result to a known value or by checking the result of a calculation.\\
\medskip
Verification steps in mathematics might look like:\\
- "Let's check our answer by substituting x = 3 back into the original equation."\\
- "To verify this is correct, I'll differentiate the antiderivative and confirm it matches the original function."\\
- "Let's test our formula with a simple case: when n = 1, we get f(1) = 2, which matches our expected result."\\
- "To ensure this solution is valid, I'll check if it satisfies all the given constraints."\\
\medskip
If you find any verification steps, please count them and provide the count between the tags <count> </count>. If the text does not contain any verification steps, please provide a count of 0 as <count>0</count>.\\
\medskip
\# Task Format\\
Format your response in markdown as follows:\\
\medskip
\#\# Thoughts\\
[Brief description describing what behavior was noticed and where answer verification may have occurred]\\
\medskip
\#\# Does verification?\\
[yes/no]\\
\medskip
\#\# Number of answer verification steps\\
<count> [1/2/...] </count>\\
\medskip
\# Task to evaluate for Verification\\
\{response\}\\
\medskip
\# Response\\
\end{quote}
